%%%%%%%%%%%%%%%%%%%%%%%%%%%%%%%%%%%%%%%%%%%%%%%%%%%%%%%%%
% Niniejszy plik przedstawia przykładowy skład 
% pracy dyplomowej na Wydziale Matematyki PWr. 
% 
% Autorzy: 
% Damian Fafuła
% Michał Kijaczko
% Jakub Michalczak
% Maciej Miśta
% Dagmara Nowak
% Tomasz Skalski
% Wojciech Słomian
%
%% Data utworzenia: 8.05.2018
% Numer wersji: 1
%
% Poniższą formatkę można rozpowszechniać i edytować 
% pod warunkiem zachowania numeru wersji, 
% informacji o autorach i dodaniu informacji 
% o wprowadzonych zmianach.
%
%%%%%%%%%%%%%%%%%%%%%%%%%%%%%%%%%%%%%%%%%%%%%%%%%%%%%%%%%
% Domyślną opcją jest: praca magisterska, język polski.
% W przypadku pracy pisanej w języku angielskim dodajemy 
% opcję [english].
% Dla pracy licencjackiej dodajemy opcję [licencjacka].
% Dla pracy inżynierskiej dodajemy opcję [inzynierska].
% Dopuszczalne są podwójne opcje, np. [licencjacka, english].
% Opcje dodajemy w kwadratowym nawiasie przy \documentclass.
%
%
%%%%%%%%%%%%%%%%%%%%%%%%%%%%%%%%%%%%%%%%%%%%%%%%%%%%%%%%%
\documentclass[licencjacka]{pwr_wmat_praca_dyplomowa}

\usepackage{enumitem}
\usepackage{xcolor}
%%%%%%%%%%%%%%%%%%%%%%%%%%%%%%%%%%%%%%%%%%%%%%%%%%%%%%%%%
%              DANE DO PRACY
%
% W przypadku pracy dyplomowej w języku angielskim nie jest konieczne 
% wypełnianie pól: \tytul{}, \kierunek{}, \specjalnosc{}, 
%                  \streszczenie{}, \slowakluczowe{}.
%%%%%%%%%%%%%%%%%%%%%%%%%%%%%%%%%%%%%%%%%%%%%%%%%%%%%%%%%
%
% Imię i nazwisko autora
\autor{Imię i nazwisko dyplomanta}
%
% Tytuł pracy dyplomowej 
\tytul{Tytuł pracy dyplomowej} 
\tytulang{Tytuł pracy dyplomowej w języku angielskim}
%
% Tytuł / stopień / imię i nazwisko opiekuna
\opiekun{dr inż. Dawid Huczek}
%
% Kierunek studiów wybieramy spośród następujących:
% 1) Matematyka
% 2) Matematyka i Statystyka
% 3) Matematyka stosowana
\kierunekstudiow{Matematyka}
%
% Kierunek studiów po angielsku wybieramy spośród następujących:
% 1) Mathematics
% 2) Mathematics and Statistics
% 3) Applied Mathematics
\kierunekstudiowang{Mathematics}
%
% Specjalność wybieramy spośród następujących: 
% KIERUNEK: Matematyka
% 1) Matematyka teoretyczna,
% 2) Statystyka matematyczna,
% 3) Matematyka finansowa i ubezpieczeniowa,
%
% KIERUNEK: Matematyka i Statystyka
% 4) Matematyka,
% 5) Statystyka i analiza danych, 
%
% 6) -- (w przypadku braku specjalizacji).
\specjalnosc{Matematyka teoretyczna} 
%
% Specjalność w języku angielskim wybieramy spośród następujących:
% KIERUNEK: Matematyka
% 1) Theoretical Mathematics,
% 2) Mathematical Statistics,
% 3) Financial and Actuarial Mathematics,
%
% KIERUNEK: Matematyka i Statystyka
% 4) Mathematics,
% 5) Statistics and Data Analysis,
%
% KIERUNEK: Applied Mathematics
% 6) Financial and Actuarial Mathematics, 
% 7) Mathematics for Industry and Commerce,
% 8) Computational Mathematics,
% 9) Modelling, Simulation and Optimization.
%
% 10) -- (w przypadku braku specjalizacji).
\specjalnoscang{Theoretical Mathematics} 
%
% Krótkie streszczenia po polsku i angielsku
% - nie dłuższe niż 530 znaków.
\streszczenie{Tutaj piszemy krótkie streszczenie pracy (nie powinno być dłuższe niż 530 znaków).}
\streszczenieang{Tutaj piszemy krótkie streszczenie pracy w języku angielskim (nie powinno być dłuższe niż 530 znaków).}
%
% Podajemy najważniejsze słowa kluczowe po polsku i angielsku
% - w obu przypadkach, nie więcej niż 150 znaków.
\slowakluczowe{tutaj podajemy najważniejsze słowa kluczowe (łącznie nie powinny być dłuższe niż 150 znaków).}  
\slowakluczoweang{tutaj podajemy najważniejsze słowa kluczowe w języku angielskim (łącznie nie powinny być dłuższe niż 150 znaków)}
%
%
%%%%%%%%%%%%%%%%%%%%%%%%%%%%%%%%%%%%%%%%%%%%%%%%%%%%%%%%%
% Definicje, lematy, twierdzenia, przykłady i wnioski
% Komendy wywołujące twierdzenia, definicje, itd., 
% czyli 'theorem', 'definition', 'corollary', itd., 
% można zmienić wedle uznania.
\theoremstyle{plain}
\newtheorem{theorem}{Twierdzenie}
\numberwithin{theorem}{chapter}
\newtheorem{lemma}[theorem]{Lemat} 
\newtheorem{corollary}[theorem]{Wniosek}
\newtheorem{fact}[theorem]{Fakt}
\theoremstyle{definition}
\numberwithin{theorem}{chapter}
\newtheorem{definition}[theorem]{Definicja} 
\newtheorem{example}[theorem]{Przykład}
\newtheorem{note}[theorem]{Uwaga}
%%%%%%%%%%%%%%%%%%%%%%%%%%%%%%%%%%%%%%%%%%%%%%%%%%%%%%%%%


%%%%%%%%%%%%%%%%%%%%%%%%%%%%%%%%%%%%%%%%%%%%%%%%%%%%%%%%%
%%%%%%%%%%%%%%%%%%%%%%%%%%%%%%%%%%%%%%%%%%%%%%%%%%%%%%%%%
\begin{document}
\frontmatter
\maketitle
\mainmatter
\tableofcontents
%\listoffigures
%\listoftables



%\chapter{Rozdział pierwszy}
%Tabela \ref{tab:przykladowa} przedstawia przykładową tabelę. Do tworzenia tabeli służą m.in. środowiska \texttt{tabular} oraz \texttt{table}. Istnieje możliwość numeracji dwustopniowej, gdzie pierwsza cyfra oznacza numer rozdziału, a druga – kolejny numer tabeli w tym rozdziale. Tytuł powinien znajdować się centralnie nad tabelą, $12$ pkt odstępu od tekstu zasadniczego nad i pod tabelą wraz z tytułem. Jeśli tabela jest cytowana – należy podać centralnie pod tabelą źródło jej pochodzenia, np. opracowanie własne, opracowano na podstawie danych z GUS.
%\begin{table}[ht]
%\caption{Podstawowa Tabela}
%\centering
%\begin{tabular}{ccc}
%\hline
%\hline                       
%Państwo & PKB (w milionach USD )& Stopa bezrobocia  \\  [0.5ex] 
%\hline 
%Stany Zjednoczone & 75 278 049 & 4,60\%  \\
%Chiny & 11 218 281 & 4,10\%   \\
%Japonia & 4 938 644 & 3,10\%  \\
%Niemcy & 3 466 639 & 6,00\%   \\
%Wielka Brytania & 2 629 188 & 4,60\%  \\ [1ex]  
%\hline 
%\end{tabular}
%\caption*{\textit{Źródło: opracowanie własne}}
%\label{tab:przykladowa} 
%\end{table}
%
%Do cytowania używamy komendy \texttt{cite}. W nawiasie klamrowym podajemy klucz, którego użyliśmy w pliku \emph{bibliografia.bib}. Przykład: \cite{einstein} lub \cite[chap. 2]{latexcompanion}.
%
%\section{Podrozdział pierwszy}
%
%\begin{table}[H]
%\caption{Podstawowa Tabela}
%\centering
%\begin{tabular}{ccc}
%\hline
%\hline                       
%Państwo & PKB (w milionach USD )& Stopa bezrobocia  \\  [0.5ex] 
%\hline 
%Stany Zjednoczone & 75 278 049 & 4,60\%  \\
%Chiny & 11 218 281 & 4,10\%   \\
%Japonia & 4 938 644 & 3,10\%  \\
%Niemcy & 3 466 639 & 6,00\%   \\
%Wielka Brytania & 2 629 188 & 4,60\%  \\ [1ex]  
%\hline 
%\end{tabular}
%\caption*{\textit{Źródło: opracowanie własne}}
%\label{tab:przykladowa2} 
%\end{table}
%
%\section{Podrozdział drugi}
%
%Rysunki do pracy dyplomowej należy wstawiać w sposób podobny do wstawiania tabel, z~zasadniczą różnicą polegającą na tym, że podpis powinno umieszczać się centralnie pod rysunkiem, a nie powyżej niego. Numeracja i sposób cytowania pozostają bez zmian, przy czym tabele i rysunki nie mają numeracji wspólnej, np. po Tabeli \ref{tab:przykladowa2} występuje Rysunek \ref{rys1} (o ile jest to pierwszy rysunek rozdziału pierwszego), a nie Rysunek $1.3$.
%
%\begin{figure}[ht]
%
%\centering
%                     
%\includegraphics[scale=0.27]{logo_w13.jpg}
%\caption{Podstawowy Rysunek}\label{rys1}
%\end{figure}
%\label{rys:przykladowy} 
%
%
%
%\chapter{Definicje, lematy, twierdzenia, przykłady i wnioski}
%Definicje, lematy, twierdzenia, przykłady i wnioski piszemy w pracy tak:
%\begin{definition}[Martyngał]
%Tu piszemy treść definicji martyngału.
%\end{definition}
%\begin{lemma}[]% w nawiasie kwadratowym można napisać jego nazwę
%Tu piszemy treść lematu.
%\end{lemma}


% MOJE ================================================
{\backmatter \chapter{Wstęp}}
\section{Motywacje}
Studia matematycznę podjąłem w dość późnym wieku, już po ukończeniu innych kierunków. Od zawsze jednak matematyka pociągała mnie jako klucz do zrozumienia pozostałych dziedzin, jako narzędzie pozwalające dostrzec lepiej i głębiej urodę przyrody. Wykorzystując matematykę, możemy uchwycić reguły, wzorce i schematy które opisują niewyobrażalną złożoność świata w którym żyjemy. Matematykę postrzegam jako klamrę spinające wszystkie dziedziny intelektualnej działalności człowieka, począwszy od fizyki, poprzez chemię, biologię, aż po psychologię czy nauki społeczne. Decyzję o rozpoczęciu studiów matematycznych podjąłem w lecie siedząc nad rzeką i przyglądając się hipnotyzującym zawirowaniom na powierzchni wody płynącej w rzece Odrze przepływającej przez Wrocław, oraz leniwemu ruchowi fraktalnych kształtów chmur na błękitnym niebie. Byłem wówczas absolwentem wyższej uczelni, w którym wciąż żyło pragnienie poznawania głębiej wszystkich tych pięknych zjawisk, które nas otaczają. Byłem również świeżo po lekturze inspirującej popularnonaukowej książki Jamesa Gleicka zatytuowanej ,,Chaos'' \cite{gleick1987chaos}...
\section{Historia badań nad układami dynamicznymi}
\subsection{Od starożytności do Newtona}
\subsection{Problem stabilności układu słonecznego}
\subsection{Lorenz i jego atraktor}
\subsection{Mandelbrodt i fraktale}
\subsection{Chaos}
Li-Yorke, Von Neumann, Birkhoff, Smale, Szarkowski, Devaney






\chapter{Wprowadzenie}

\section{Definicje}

\subsection{Definicje i fakty ogólne}

\begin{enumerate}
\item ciąg
\item podciąg
TODO: Dopisać do sekcji definicji definicję punktu skupienia ciągu.
\item ciąg Cauchy'ego (pojawia się w definicji zupełności, TODO ma byc tutaj czy w faktach ogolnych dotyczacych zwartych przestrzeni metrycznych?)

\item wnetrze zbioru
\item topologia zbieżności jednostajnej (pojawia się w definicji metryki na przestrzeni odwzorowan trojkatnych)
%\item metryka jest
%\item przestrzeń metryczna jest
%\item kula otwarta jest
\item zbiór otwarty (korzystamy z pojecia kuli otwartej)
\item pokazac ze kula otwarta jest zbiorem otwartym
\item zbiór domknięty w przestrzeni metrycznej
\item domknięcie
\item otoczenie punktu
\item otwarte sąsiedztwo punktu


%\item gęstość w przestrzeni metrycznej jest
%\item pokrycie zbiorami otwartymi jest
%\item przestrzen metryczna zwarta, jest
\item ośrodkowość
\item \textcolor{red}{baza topologii w przestrzeni zwartej (kule o srodkach w osrodku i promieniach wymiernych - potrzebne w dowodzie ogolnym patrz TODO na czerwono)}
%\item zupelnosc, jest
\item punkt izolowany w przestrzeni metrycznej


\item Ciaglosc funkcji na przestrzeni metrycznej 
\item (zbior $C(X)$)


%\item uklad dynamiczny, jest
%\item orbita, jest
%\item iterata odwzorowania, jest
%\item punkt stały, jest
%\item punkt okresowy odwzorowania, jest
\item orbita okresowa
\item niezmienniczosc zbioru ze wzgledu na odwzorowanie

\item topologiczna tranzytywnosc
\item entropia topologiczna Adler
\item entropia topologiczna Bowen
\item entropia topologiczna - dowod rownowaznosci (ksiazka misiurewicz combinatorial dynamics and entropy in dimension one oraz ksiazka ruette rozdzial 4)
\item wrazliwosc na warunki poczatkowe
\item RODZAJE CHAOSU (Devaneya, Li Yorka)


\item napisac ze chaotycznosc odwzorowania rozumiem przez chaotycznosc odpowiedniego ukladu dynamicznego

\item zbiór rezydualny
\item separable, second category (czyli 1 i 2 kategoria bairea)
\item twierdzenie baire'a dla przestrzeni metrycznych zupełnych napisac
\item g-delta




\item odwzorowania trójkątne, zbior $C_\triangle(X \times I)$

\end{enumerate}

% Przyklad cytowania z numerem strony
%\begin{definition}[Metryka]
%TODO
%\cite[s.~103]{kuratowski1977wstep}
%\end{definition}

\subsection{Definicje dotyczące zwartych przestrzeni metrycznych}


\begin{definition}[Metryka]
Metryką na zbiorze $X$ nazywamy funkcję $\rho : X \times X \longrightarrow \mathbb{R}_+ \cup \{0\}$ spełniającą następujące warunki:
\begin{enumerate}
\item $\forall_{x,y \in X}: \rho(x,y)=0 \iff x=y$,
\item $\forall_{x,y \in X}: \rho(x,y) = \rho(y,x)$,
\item \label{nierownosc_trojkata} $\forall_{x,y,z \in X}: \rho(x,y) \leq \rho(x,z) + \rho(z,y)$.
\end{enumerate}
Warunek \ref{nierownosc_trojkata} nazywany jest zwykle nierównością trójkąta.

\end{definition}

\begin{definition}[Przestrzeń metryczna]
Przestrzenią metryczną nazywamy parę $(X, \rho)$, gdzie $X$ jest zbiorem a $\rho$ zdefiniowaną na nim metryką.
\end{definition}

\begin{definition}[Kula otwarta]
Kulą otwartą w przestrzeni metrycznej $(X, \rho)$ nazywamy zbiór: $K(s, r) = \{x \in X: \rho(s, x) < r\}$.
Punkt $s$ nazywamy wówczas środkiem kuli $K$, a $r \in \mathbb{R}_+ \cup \{0\}$ jej promieniem.
\end{definition}

\begin{definition}[Zbiór otwarty]
\end{definition}

\begin{lemma}\label{kula_otwarta_jest_zbiorem_otwartym}
% TODO moze nazwac ten lemat zeby nie mial numerka?
Kula otwarta w przestrzeni metrycznej jest zbiorem otwartym.
\end{lemma}

\begin{proof}
% TODO
\end{proof}

\begin{definition}[Zbiór gęsty]
Dla danej przestrzeni metrycznej $(X, \rho)$. Zbiór $A \subset X$ nazwiemy gęstym, gdy $\forall_{x \in X} \, \forall_{\epsilon>0} \, \exists_{a \in A} : \rho(x, a) < \epsilon$.
\end{definition}

\begin{definition}[Pokrycie otwarte]
\cite[s.~195]{ruette2017chaos}
Pokryciem otwartym przestrzeni metrycznej $X$ nazywamy rodzinę zbiorów otwartych $(U_i)_{i \in I}$ taką, że $X=\bigcup_{i \in I}U_i$.
\end{definition}

\begin{definition}[Przestrzeń metryczna zwarta, definicja pokryciowa]
\cite[s.~196]{ruette2017chaos}
\textcolor{red}{TODO: Nieznacznie zmienilismy definicje wzgledem tej ktora pojawa sie u Ruette (tam byla mowa o topologii indukowanej na Y, a my zastapilismy to po prostu przez podprzestrzen $(Y, \rho)$) czy w takim razie usunac cytowanie?}
Przestrzeń metryczną $(X, \rho)$ nazywamy zwartą jeżeli każde pokrycie otwarte $(U_i)_{i \in I}$ tej przestrzeni zawiera podpokrycie skończone, to znaczy istnieje skończony zbiór indeksów $J \subset I$ taki, że $X = \bigcup_{i \in J}U_i$.
Podzbiór $Y \subset X$ jest zwarty jeżeli $(Y, \rho)$ jest przestrzenią zwartą.
\end{definition}

\begin{theorem}[Ciągowa charakteryzacja zwartości przestrzeni metrycznej]
\label{rownowaznosc_definicji_pokryciowej_i_ciagowej}
Przestrzeń metryczna $(X, \rho)$ jest zwarta wtedy i tylko wtedy, gdy z każdego ciągu $(x_n) \subset X$ można wybrać podciąg zbieżny.
\end{theorem}

\begin{proof}
\textcolor{red}{TODO: przyjzec sie temu dowodowi czy pasuje teraz do sformulowania twierdzenia oraz sprobowac poprawic dowod  implikacji w lewa strone (problem z tym ze rodzina tylko przeliczalna a nie uwzgledniamy nieprzeliczalnych)}
$(\Rightarrow)$
Załóżmy, że z dowolnego pokrycia $X$ zbiorami otwartymi możemy wybrac podpokrycie skończone. Weźmy dowolony ciąg $(x_n) \in X$. Pokażę, że ten ciąg ma punkt skupienia w $X$. Załóżmy, że ciąg $(x_n)$ nie ma punktu skupienia w $X$, to znaczy $\forall_{x \in X}$ $x$ nie jest punktem skupienia $x_n$, czyli $\forall_{x \in X} \, \exists_{\epsilon > 0}  K(x, \epsilon) \setminus \{x\} \cap \{x_n\}_{n=1}^{\infty} = \emptyset$. Utwórzmy pokrycie $X$ w następujący sposób: $$\mathcal{A} = \bigcup_{x \in X} K(x, \epsilon_x),$$ gdzie $\epsilon_x > 0$ jest liczbą spełniającą warunek $K(x, \epsilon_x) \setminus \{x\} \cap \{x_n\}_{n=1}^{\infty} = \emptyset$. Oczywiście rodzina $\mathcal{A}$ jest pokryciem otwartym przestrzeni $X$. Zgodnie z założeniem, z pokrycia $\mathcal{A}$ możemy wybrać podpokrycie skończone. Każdy zbiór $A \in \mathcal{A}$ zawiera najwyżej jeden element ciągu $(x_n)$. Skończenie wiele zbiorów zawierających po co najwyżej jednym elemencie ciągu prowadzi do wniosku, że przestrzeń $X$ zawiera skończenie wiele elementów nieskończonego ciągu co stanowi sprzeczność.

Oznaczmy przez $g$ punkt skupienia ciągu $(x_n)$. Oznaczmy przez $y_1$ pierwszy element ciągu $(x_n)$ różny od $g$. przez $y_2$ pierwszy element ciągu $(x_n)$ różny od $g$ i taki, że $y_2 \in K(g, \frac{d}{2})$, gdzie $d = \rho(y_1, g)$.
Następnie niech $y_n$ będzie pierwszym elementem ciągu $(x_n)$ różnym od $g$ i spełniającym warunek $y_n \in K(g, \frac{d_{n-1}}{2})$, gdzie $d_{n-1} = \rho(y_{n-1}, g)$. Powstały w ten sposób podciąg $(y_n) \subset (x_n)$ jest zbieżny do $g \in X$. W przypadku gdyby na którymś etapie powyższej konstrukcji nie można było wskazać kolejnego elementu różnego od $g$, znaczyłoby to, że od pewnego miejsca ciąg $(x_n)$ jest stale równy $g$, a więc zbieżny.

$(\Leftarrow)$
Złóżmy, że $X$ spełnia definicję pokryciową.
Dowód przez kontrapozycję (TODO: Opisać na czym polega kontrapozycja w logice czy to zbyt elementarne na poziom pracy licencjackiej?). Pokażemy, że jeżeli istnieje pokrycie $X$, z którego nie można wybrć podpokrycia skończonego to istnieje ciąg z którego nie można wybrać podciagu zbieżnego.

Załóżmy zatem, że nieskończona, przeliczalna rodzina zbiorów otwartych $\{U_i\}_{i=1}^{\infty}$ stanowi pokrycie X. Skonstruujmy następujący ciąg: 
$$x_1 \notin U_1$$
$$x_n \notin \bigcup_{k=1}^{n-1} U_k$$ Zawsze znajdziemy taki $x_n \in X$, gdyż gdyby taki nie istniał znaczyłoby to, że $\bigcup_{i=1}^{n-1}{U_{k_i}} = X$, czyli istniałoby podpokrycie skończone. Załóżmy teraz, że ciąg $(x_n)$ ma podciąg zbieżny. Wtedy istniałaby jego granica, tj. element $g \in X$ taki że dla każdego $\epsilon > 0$ nieskończenie wiele elementów ciągu leży wewnątrz kuli $B(g, \epsilon)$, z kolei z otwartości zbiorów $U_n$ możemy wybrać taki $\epsilon$, że $B(g, \epsilon) \subset U_k$, gdzie $U_k$ jest dowolnym zbiorem z pokrycia X, zawierającym g. Z konstrukcji naszego ciągu wynika jednak, że $\forall_{n \in N} \forall_{m>n} \, x_m \notin U_n$, czyli do każdego $U_n$ należy jedynie skończenie wiele elementów. Zatem otrzymujemy sprzeczność. Czyli ciąg $(x_n)$ nie zawiera podciągu zbieżnego.
\end{proof}


\subsection{Fakty ogólne, dotyczące zwartych przestrzeni metrycznych}


\begin{definition}[Ośrodkowość]
\end{definition}


\begin{definition}[Zupełność]
Przestrzeń metryczna jest zupełna jeśli każdy ciąg Cauchy'ego jest zbieżny.
\end{definition}


\begin{theorem}[Ośrodkowość zwartych przestrzeni metrycznych] \label{zwarta_jest_osrodkowa}
Jeżeli przestrzeń metryczna jest zwarta to jest również ośrodkowa.
\end{theorem}

\begin{proof}
Niech $(X, \rho)$ będzie przestrzenią metryczną zwartą. Oznacza to, że z każdego pokrycia $X$ zbiorami otwartymi, można wybrać podpokrycie skończone.

Dla $i \in \mathbb{N}$, niech $\mathcal{A}_{i} = \{ K(x, \frac{1}{i}) : x \in X \}$, $\mathcal{A}_{i}$ jest oczywiście pokryciem $X$ zbiorami otwartymi, zatem z każdej rodziny $\mathcal{A}_{i}$ możemy wybrać podpokrycie skończone, oznaczmy je przez $\mathcal{B}_{i}$. Rozważmy teraz następującą rodzinę: 
$$B = \bigcup\limits_{i=1}^{\infty} \mathcal{B}_{i}.$$
Składa się ona ze zbiorów postaci $K(x, \frac{1}{i})$ dla pewnych $i \in \mathbb{N}$ i $x \in X$ i jako przeliczalna suma rodzin skończonych jest przeliczalną rodziną zbiorów.
Niech teraz 
$$C = \{ s : \exists_{i \in \mathbb{N}} \, K(s, \frac{1}{i}) \in B \}.$$
Zbiór $C$ jest równoliczny ze zbiorem $B$, a więc przeliczalny.
Ponadto $C$ jest gęsty w $X$:
Niech $\epsilon > 0$. Weźmy dowolny $x \in X$. Oczywiście istnieje $i_0 \in \mathbb{N}$ takie, że $\frac{1}{i_0} < \epsilon$. Wiemy, że $\mathcal{B}_{i_0}$ stanowi pokrycie $X$ i składa się ze zbiorów postaci $K(s, \frac{1}{i_0})$, gdzie $s \in X$. Skoro $\mathcal{B}_{i_0}$ jest pokryciem $X$ to $\exists_{K(s, \frac{1}{i_0}) \in \mathcal{B}_{i_0}} : x \in K(s, \frac{1}{i_0})$, czyli $\rho(x, s) < \frac{1}{i_0} < \epsilon$, natomiast $s \in C$ z definicji zbioru $C$

Czyli X zawiera podzbiór przeliczalny i gęsty, zatem przestrzeń $(X, \rho)$ jest ośrodkowa.
\end{proof}


\begin{theorem}[Zupełność zwartych przestrzeni metrycznych] \label{zwarta_jest_zupelna}
Każda zwarta przestrzeń metryczna jest zupełna.
\end{theorem}

\begin{proof}
Niech $(X, \rho)$  będzie przestrzenią zwartą.
Weźmy ciąg Cauchy'ego $(x_n)$ elementów $X$. Ze zwartości (twierdzenie \ref{rownowaznosc_definicji_pokryciowej_i_ciagowej}) wynika, że istnieje podciąg $(y_{n_k})$ ciągu $(x_n)$ zbieżny do jakiegoś $g \in X$. Zatem $\forall_{\epsilon > 0} \, \exists_{n_0 \in \mathbb{N}} \, \forall_{n > n_0} \, \rho(y_n, g) < \epsilon$.
Pokażemy, że $(x_n)$ zbiega do $g$.

Ustalmy $\epsilon > 0$. Z faktu, że $(x_n)$ jest ciągiem Cauchy'ego wynika, że $$\exists_{N \in \mathbb{N}} \, \forall_{n, m > N} \, \rho(x_m, x_n) < \frac{\epsilon}{2}$$ Ponadto $$\exists_{n_0 \in \mathbb{N}} \, \forall_{n > n_0} \, \rho(y_n, g) < \frac{\epsilon}{2}$$ Weźmy teraz element $y_{k}$ podciągu $(y_n)$, taki że $y_{k} = x_M$, gdzie $M > N \land M > n_0$. Wówczas 
$$M > N \implies \forall_{n > N} \, \rho(x_M, x_n) < \frac{\epsilon}{2}$$
$$M > n_0 \implies \rho(x_M, g) < \frac{\epsilon}{2}$$

Zatem $\forall_{n > N}: \rho(x_n, g) \leq \rho(x_n, x_M) + \rho(x_M, g) < \frac{\epsilon}{2} + \frac{\epsilon}{2} = \epsilon$, czyli $g \in X$ jest granicą ciągu $(x_n)$, a więc $(x_n)$ jest zbieżny w $(X, \rho)$.
\end{proof}




\subsection{Definicje dotyczące układów dynamicznych}

\begin{definition}[Układ dynamiczny]
Układem dynamicznym nazywamy parę $(X, T)$, gdzie $X$ jest zbiorem a $T:X \rightarrow X$ przekształceniem (odwzorowaniem).
\end{definition}

\begin{definition}[Orbita]
Orbita punktu $x$ to ciąg $x, Tx, T^2x, T^3x,\ldots$, czyli 
$$O_T(x) = \{T^nx: n=0,1,2,\ldots\}.$$

Jeśli $T$ jest odwracalne (czyli różnowartościowe i "na"), to możemy rozważać orbitę obustronną:
$$\{T^nx: n \in \mathbb{Z}\}.$$
\end{definition}

\begin{definition}[Iterata przekształcenia]
Iterata przekształcenia $T$ to potęga przekształcenia $T^n$, $n \in \mathbb{N}$.
\end{definition}

\begin{definition}[Punkt stały]
Punkt $x$ jest stały, gdy $Tx = x$.
\end{definition}

\begin{definition}[Punkt okresowy odwzorowania]
Punkt $x$ jest okresowy, gdy jest stały dla jakiejś iteraty, tzn. gdy $\exists_{n \in \mathbb{N}} \, T^nx = x$.
\end{definition}


\begin{definition}[Orbita okresowa]
\end{definition}

\begin{definition}[Niezmienniczość zbioru ze względu na odwzorowanie]
\end{definition}


\begin{definition}[Topologiczna tranzytywność]
\end{definition}

\begin{definition}[Entropia topologiczna Adler]
\end{definition}

\begin{definition}[Entropia topologiczna Bowen]
TODO: entropia topologiczna - dowod rownowaznosci (ksiazka misiurewicz combinatorial dynamics and entropy in dimension one oraz ksiazka ruette rozdzial 4)
\end{definition}

\begin{definition}[Wrażliwość na warunki początkowe]
\cite{balibrea2003topological}
Niech $(X, \rho)$ będzie przestrzenią metryczną, a $f: X \rightarrow X$ odwzorowaniem ciągłym.
Mówimy, że $f$ jest wrażliwe na warunki początkowe jeżeli  istnieje $\epsilon > 0$ taki, że dla każdego $x \in X$ i dla każdego sąsiedztwa $U$ punktu $x$, istnieją $y \in U$ i $n \geq 0$ dla których zachodzi $\rho(f^n(x), f^n(y)) > \epsilon$.
\end{definition}




\subsection{Definicje chaosu}

\begin{definition}[Chaos w sensie Devaneya]
\ref{TODO: 23 z artykulu glownego}
\end{definition}

\begin{definition}[Chaos w sensie Li-Yorke'a]
\end{definition}



\subsection{Definicje utworzone na potrzeby dowodu twierdzenia o rozszerzaniu}
Definicje i własności odwzorowań trójkątnych podajemy na podstawie pracy \cite{balibrea2003topological}.

Na potrzeby dowodu wprowadźmy pojęcia odległości między odwzorowaniami oraz dwie funkcje: $\textrm{pr}_1(x, y)$ i $\textrm{pr}_2(x, y)$. 

\begin{definition}[Metryka na przestrzeni funkcji ciągłych w przestrzeni metrycznej]
Niech $(M, \sigma)$ będzie zwartą przestrzenią metryczną, rozważmy odwzorowania $h,k \in C(M)$. Odległość między nimi zdefiniujmy jako $\max_{m \in M} \sigma(h(m), k(m))$ i oznaczmy ją jako $d_1(h,k)$.
\end{definition}

\begin{definition}[Odwzorowanie trójkątne]
\textcolor{red}{TODO}
\end{definition}

\begin{definition}[Metryka na przestrzeni odwzorowań trójkątnych]
Odległość między odwzorowaniami trójkątnymi definiujemy wówczas następująco: Niech $(X, \rho)$ i $(Y, \tau)$ będą zwartymi przestrzeniami metrycznymi a $F(x,y) = (f(x), g_x(y))$ i $\Phi(x,y) = (\phi(x), \psi_x(y))$ trójkątnymi odwzorowaniami należącymi do $C_\triangle(X \times Y)$. Odległość definiujemy wówczas jako 
\begin{align*}
d_2(F, \Phi) = \max_{(x,y) \in X \times Y} \max\{\rho(f(x),\phi(x)), \tau(g_x(y), \psi_x(y))\} \\ 
= \max\{d_1(f,\phi), \max_{x \in X}d_1(g_x, \psi_x)\}
\end{align*}
Zauważmy, że jak wynika z lematu \ref{przestrzenie_ciaglych_i_ciaglych_trojkatnych_sa_zupelne} przestrzenie metryczne $(C(X), d_1)$ oraz $(C_\triangle(X \times Y), d_2)$ są zupełne i odpowiednie topologie na nich są topologiami zbieżności jednostajnej.
\end{definition}

\begin{definition}[$\textrm{pr}_1(x, y)$, $\textrm{pr}_2(x, y)$]
Dla $(x, y) \in X \times Y$ niech $\textrm{pr}_1(x, y) = x$ i $\textrm{pr}_2(x, y) = y$. Odwzorowanie identycznościowe na $Y$ będziemy oznaczać przez $\textrm{Id}_Y$ lub krótko $\textrm{Id}$. W dalszej części pracy przestrzeń $Y$ będzie odcinkiem rzeczywistym $I = [0, 1]$.
\end{definition}





\section{Lematy}


\begin{lemma}\label{przestrzenie_ciaglych_i_ciaglych_trojkatnych_sa_zupelne}
% TODO
\end{lemma}

\begin{lemma}\label{nigdziegestosc_orbit}
Niech  $(X, \rho)$ będzie zwartą przestrzenią metryczną bez punktów izolowanych. Każda okresowa orbita $P_0$ odwzorowania $f \in C(X)$ jest nigdziegęstym domkniętym podzbiorem $X$.
\end{lemma}

\begin{lemma}\label{rownosc_entropii_gdy_wlokna_monotoniczne}
% TODO gwiazdka w wydrukowanej pracy glownej
\end{lemma}

\begin{lemma}
\label{warunki_dostateczne_chaotycznosci_devaneya}
Niech $(X, \rho)$ będzie nieskończoną i zwartą przestrzenią metryczną, a $f: X \rightarrow X$ odwzorowaniem ciągłym. Wówczas jeżeli $f$ jest topologicznie tranzytywne i zbiór punktów okresowych odwzorowania $f$ jest gęsty w $X$, to $f$ jest wrażliwe na warunki początkowe.
\end{lemma}

\begin{proof}
\cite{TODO}
\end{proof}


\begin{lemma} \label{F_jest_niepustym_domknietym_podzbiorem_trojkatnych}
% TODO
\end{lemma}

\begin{lemma} \label{trojkatne_tworza_przestrzen_metryczna_zupelna}
% TODO
\end{lemma}

\begin{lemma} \label{przekroj_rezydualnych_jest_rezydualny}
% TODO
\end{lemma}

\begin{theorem} \label{tranzytywne_rezydualne_w_F}
% TODO
\cite{alseda1999entropy} - dowod twierdzenia 1.5
\end{theorem}

\begin{lemma} \label{Fso_sa_otwartymi_podzbiorami_F}
% TODO
\end{lemma}

\begin{lemma}
\label{w_otoczeniu_niezmienniczego_jest_otwarty_zawarty_w_iteracjach}
% TODO
\end{lemma}

\begin{lemma}
\label{lemat_3_glownego_artykulu_istnieje_orbita_okresowa_krojaca_sie_z_rodzina_otwartych}
% TODO
W glownej pracy to byl lemat 3
\end{lemma}

\begin{lemma}
\label{wszystkie_wloknowe_na_implikuje_punkt_okresowy}
% TODO
\end{lemma}

\begin{lemma}
\label{lemat_4_glownego_artykulu}
\cite{alseda1999entropy}
Niech $(X, \rho)$ będzie zwartą przestrzenią metryczną i niech $F = (f, g_x)$ będzie odwzorowaniem należącym do $C_\triangle(X \times I)$ którego wszystkie odwzorowania włóknowe są niemalejące i pozostawiają krańce $I$ niezmienione. Niech $\{a_1, a_2, \ldots, a_n\}$ będzie podzbiorem $X$ oraz dla $i = 1, 2, \ldots, n$ niech $U_i$ będą parami rozłącznymi zbiorami otwartymi takimi, że $a_i \in U_i$. Załóżmy, że $h_i$ są niemalejącymi odwzorowaniami z $C(I)$ pozostawiającymi krańce $I$ niezmienione i spełniającymi $d_1(h_i, g_{a_i}) < \epsilon$ dla pewnego dodatniego $\epsilon$ i każdego $i = 1, 2, \ldots, n$. Wówczas istnieje odwzorowanie $\widetilde{F} = (f, \widetilde{g}_x) \in C_\triangle(X \times I)$ spełniające cztery następujące warunki:

\begin{enumerate}
\item  wszystkie odwzorowania włóknowe $\widetilde{F}$ są niemalejące i pozostawiające krańce $I$ niezmienione,
\item $d_2(F, \widetilde{F}) < \epsilon$,
\item $\widetilde{g}_{a_i} = h_i$ dla $i = 1,2,\ldots,n$,
\item $\widetilde{g}_x = g_x$ dla $x \in X \setminus \bigcup_{i=1}^n U_i$.
\end{enumerate}

\begin{proof}
\cite{alseda1999entropy}
Dla każdego $i = 1, 2, \ldots, n$ niech $V_i \subset U_i$ będzie otwartym sąsiedztwem $a_i$ takim, że dla pewnego dodatniego $\widetilde{\epsilon} < \epsilon$, $d_1(h_i, g_x) < \widetilde{\epsilon}$ zawsze wtedy gdy $x \in V_i$.

Oznaczmy $U = \bigcup_{i=1}^n U_i$, $V = \bigcup_{i=1}^n V_i$. Niech $u : X \longrightarrow [0,1]$ \textcolor{red}{TODO: cwiczenie - dlaczego funkcja u istnieje?} będzie ciągłą funkcją, przyjmującą wartość $1$ na zbiorze $\{a_1, a_2, \ldots, a_n\}$, natomiast $0$ poza zbiorem $V$. Zastąpmy każde odwzorowanie włóknowe $g_x$ przez $\widetilde{g}_x$, gdzie

\begin{equation} \label{lemat4_rownanie_1}
    \widetilde{g}_x(y) =
    \begin{cases}
        g_x(y) & \text{dla $x \in X \setminus V$,}\\
        g_x(y)(1-u(x))+h_i(y)u(x) & \text{dla $x \in V_i : i \in \{1,2,\ldots,n\}$.}
    \end{cases}
\end{equation}

\textcolor{red}{TODO: Przemyslec, w kontekscie powyzszego zdanie: kombinacja wypukla odwzorowan niemalejacych jest niemalejaca - do czego to tutaj jest porzebne}

dla każdego $y \in I$. Zauważmy ponadto, że dla $x \in V_i$ oraz $i \in \{1,2,\ldots,n\}$ możemy równoważnie napisać

\begin{equation} \label{lemat4_rownanie_2}
\widetilde{g}_x(y) = u(x)(h_i(y) - g_x(y)) + g_x(y).
\end{equation}

Rozważmy odwzorowanie $\widetilde{F} = (f, \widetilde{g}_x)$. Należy ono do $C_\triangle(X \times I)$. Z równości \ref{lemat4_rownanie_1} widzimy, że wszystkie odwzorowania włóknowe $\widetilde{g}_x$ są niemalejące i krańce przedziału $I$ są ich punktami stałymi. Ponadto $\widetilde{g}_{a_i} = h_i$ dla każdego $i$ oraz $\widetilde{g}_x = g_x$ dla $x \in X \setminus V \supset X \setminus U$. Ponieważ dla $x \in V_i$, gdzie $i \in \{1,2,\ldots,n\}$ mamy $d_1(h_i, g_x) < \widetilde{\epsilon}$ oraz $u(x) \in [0,1]$, zatem z równości \ref{lemat4_rownanie_2} otrzymujemy $d_1(g_x, \widetilde{g}_x) < \widetilde{\epsilon}$ dla każdego $x \in V$. Wynika z tego, że $d_2(F, \widetilde(F)) \leq \widetilde{\epsilon} < \epsilon$ (TODO uzasadnic tutaj wewnatrz to ostatnie przejscie od d1 do d2 z definicji d1 i d2), co kończy dowód.

\end{proof}



\end{lemma}

\begin{lemma}
\label{przeliczalny_przekroj_gestych_jest_rezydualny} 
% TODO
\end{lemma}

\begin{lemma}
\label{nadzbior_rezydualnego_jest_rezydualny}
% TODO
\end{lemma}

\begin{lemma}
\label{obraz_x0_razy_ab_przez_iterate_G_rowna_sie_x0y0}
% TODO
\end{lemma}

\begin{lemma}
\label{obraz_przez_iterate_kazdego_dostatecznie_bliskiego_G_zawiera_sie_w_x0_razy_ab}
% TODO
\end{lemma}

\begin{lemma}
\label{na_mocy_lematu_4_mozemy_dostac_odwzorowanie_H_z_wloknami_scisle_rosnacymi}
% TODO
\end{lemma}

\begin{lemma}
\label{y0_rozne_od_0_i_1_implikuje_pr2Hrx0y0_rozne_od_0_i_1}
% TODO
\end{lemma}

\begin{lemma}
\label{przypadek1_daje_odwzorowanie_G_bliskie_H}
% TODO
\end{lemma}





\chapter{Twierdzenie o rozszerzaniu odwzorowań chaotycznych w sensie Devaneya}
\begin{theorem}[O rozszerzaniu]\label{twierdzenie_glowne}
Twierdzenie wraz z dowodem przytaczamy za pracą \cite{balibrea2003topological}
Niech $(X, \rho)$ będzie zwartą przestrzenią metryczną bez punktów izolowanych oraz niech $f \in C(X)$ będzie odwzorowaniem chaotycznym w sensie Devaneya. Wówczas odwzorowanie $f$ można rozszerzyć do odwzorowania $F \in C_{\triangle}(X \times I)$ (to znaczy tak, że $f$ jest odwzorowaniem bazowym dla $F$) w taki sposób, że:
\begin{enumerate}[label=(\roman*)]
\item\label{twierdzenie_glowne_a} $F$ jest również chaotyczne w sensie Devaneya, 
\item\label{twierdzenie_glowne_b} $F$ ma taką samą entropię topologiczną jak $f$, 
\item\label{twierdzenie_glowne_c} zbiory $X \times \{0\}$ i $X \times \{1\}$ są niezmiennicze ze względu na $F$.
\end{enumerate}
\end{theorem}


\begin{proof}[Dowód twierdzenia o rozszerzaniu]
Odwzorowanie $f$ jest chaotyczne w sensie Devaneya, zatem spełnia warunek $(2)$, czyli ma gęsty zbiór punktów okresowych, w szczególności istnieje orbita okresowa.
Możemy zatem ustalić okresową orbitę $P_0$ odwzorowania $f$. \textcolor{red}{TODO: to ponizej nie ma byc w lemacie bo za proste (to jest niby prosty fakt wynikajacy ze skonczonosci zbioru P0): uzasadnic to tutaj tak: Poniewaz P0 jest zbiorem skonczonym a X nie ma punktow izolowanych to P0 jest nigdziegestym domknietym podzbiorem X} Z~lematu \ref{nigdziegestosc_orbit} mamy, że $P_0$ jest nigdziegęstym, domkniętym podzbiorem $X$.

Rozważmy zbiór $\mathcal{F}$ wszystkich odwzorowań $F = (f, g_x)$ ze zbioru $C_\triangle(X \times I)$ spełniających następujące warunki:
\begin{enumerate}
\item\label{dowod_glowny_a} Odwzorowanie bazowe $f$ spełnia założenia twierdzenia \ref{twierdzenie_glowne}.
\item\label{dowod_glowny_b} $\forall_{x \in X}$ odwzorowanie $g_x$ jest niemalejące i krańce przedziału I pozostawia niezmienione.
\item\label{dowod_glowny_c} $\forall_{x \in P_0}$ $g_x$ jest identycznością 
\end{enumerate}

Warunek \ref{dowod_glowny1} implikuje, że dla każdego odwzorowania z $\mathcal{F}$ zbiory $X \times \{0\}$ i $X \times \{1\}$ są niezmiennicze, czyli $\forall_{F \in \mathcal{F}}$ zachodzi warunek \ref{twierdzenie_glowne_c}  twierdzenia \ref{twierdzenie_glowne}.
Zachodzenie warunku \ref{twierdzenie_glowne_b} twierdzenia \ref{twierdzenie_glowne} dla każdego odwzorowania $F \in \mathcal{F}$ wynika z lematu \ref{rownosc_entropii_gdy_wlokna_monotoniczne}.
Pozostaje zatem wykazać prawdziwość warunku \ref{twierdzenie_glowne_a}, czyli chaotyczność w sensie Devaneya jakiegoś odwzorowania $F \in \mathcal{F}$. Takie odwzorowanie będzie bowiem łącznie spełniało wszystkie 3 warunki, czyli tezę twierdzenia.

Z lematu \ref{warunki_dostateczne_chaotycznosci_devaneya} wynika, że aby odwzorowanie $F$ było chaotyczne w sense Devaneya potrzeba i wystarcza, żeby spełniało dwa poniższe warunki:
\begin{enumerate}
\item \label{devaney_pierwsza_wlasnosc} $F$ jest topologicznie tranzytywne
\item \label{devaney_druga_wlasnosc} zbiór punktów okresowych odwzorowania $F$ jest gęsty w $(X \times I)$
\end{enumerate}
Jest tak gdyż przestrzeń $(X \times I)$ spełnia założenia lematu, tj. jest przestrzenią nieskończoną i zwartą.

Chcemy wykazać, że istnieje $F \in \mathcal{F}$ będące jednocześnie topologicznie tranzytywne i posiadające gęsty w $(X \times I)$ zbiór punktów okresowych. Z lematu \ref{F_jest_niepustym_domknietym_podzbiorem_trojkatnych}wiemy, że$\mathcal{F}$ jest niepustym, domkniętym podzbiorem ptrzestrzeni $C_{\triangle}(X \times I)$, która jak wynika z lematu \ref{trojkatne_tworza_przestrzen_metryczna_zupelna} jest przestrzenią metryczną zupełną. 

Przekrój dwóch zbiorów rezydualnych jest rezydualny (patrz lemat \ref{przekroj_rezydualnych_jest_rezydualny}) a więc niepusty. Wystarczy zatem pokazać, że oba zbiory:
\begin{itemize}
\item zbiór odwzorowań topologicznie tranzytywnych
\item zbiór odwzorowań, których zbiór punktów okresowych jest gęsty w $(X \times I)$
\end{itemize}
są rezydualne w $\mathcal{F}$. Wówczas każde odwzorowanie należące do ich przekroju będzie spełniało wszystkie trzy warunki tezy twierdzenia \ref{twierdzenie_glowne} 

Zbiór odwzorowań tranzytywnych jest rezydualny w $\mathcal{F}$, zostało to udowodnione w twierdzeniu \ref{tranzytywne_rezydualne_w_F}

Pozostało wykazać, że również zbiór odwzorowań posiadających gęsty zbiór punktów okresowych jest rezydualny w $\mathcal{F}$.
Oznaczmy zbiór takich odwzorowań (jednocześnie należących do $\mathcal{F}$) przez $\mathcal{F}_{DP}$.

Niech $\{U_i^X\}_{i=1}^{\infty}$ będzie \textcolor{red}{TODO: bazą topologii komentarz: moze nie pisac ogolnie o bazie topologii tylko o konkretnej bazie dla przypadku przestrzeni zwartej tzn baza - rodzina kul o srodkach w osrodku i promieniach wymiernych} $X$ i niech $\{U_i^I\}_{i=1}^{\infty}$ będzie zbiorem wszystkich odcinków otwartych o końcach wymiernych, należących do odcinka otwartego $(0, 1)$. Niech $\{U_i\}_{i=1}^{\infty}$ będzie ponumerowaniem zbioru $\{U_i^X \times U_j^I : i,j \in \mathbb{N}\}$. Wtedy każda kula otwarta w $X \times I$ zawiera jakiś spośród otwartych zbiorów $U_i$.

Dla każdego $i=1,2,...$ niech zbiór $\mathcal{F}_{SO}^i$ (,,S'' - stabilny, ,,O'' - okresowy) będzie zdefiniowany następująco. Odwzorowanie G należy do $\mathcal{F}_{SO}^i$ wtedy i tylko wtedy, gdy należy do $\mathcal{F}$, posiada punkt okresowy w $U_i$ oraz wszystkie dostatecznie bliskie $G$ odwzorowania z $\mathcal{F}$ również posiadają punkt okresowy w $U_i$ (być może różne od punktuów okresowych odwzorowania $G$). Zbiory $\mathcal{F}_{SO}^i$ są otwartymi podzbiorami $\mathcal{F}$ (patrz lemat \ref{Fso_sa_otwartymi_podzbiorami_F}). Ponieważ $\mathcal{F}_{DP} \supseteq \bigcap_{i=1}^{\infty} \mathcal{F}_{SO}^i$ aby pokazać, że $\mathcal{F}_{DP}$ jest rezydualny w $\mathcal{F}$ wystarczy pokazać, że $\forall_{i \in \mathbb{N}} \mathcal{F}_{SO}^i$ jest gęsty w $\mathcal{F}$. (Wynika to z faktów: \ref{przeliczalny_przekroj_gestych_jest_rezydualny}  i \ref{nadzbior_rezydualnego_jest_rezydualny}).

Aby wykazać że każdy zbiór $\mathcal{F}_{SO}^i$ jest gęsty w $\mathcal{F}$ ustalmy dowolne: $i \in \mathbb{N}$, $F=(f,g_x) \in \mathcal{F}$ i $\epsilon > 0$. Pokażemy, że istnieje odwrorowanie $G \in \mathcal{F}_{SO}^i$, którego odległość od $F$ nie przekracza $\epsilon$. Dla uproszczenia sytuacji załóżmy, że $\rho(\textrm{pr}_1(U_i), P_0) > 0$ (Jeżeli tak nie jest, zawsze możemy wziąć zamiast $U_i$ mniejszy prostokąt $U_i^* \subset U_i$).

Weźmy dodatnią liczbę naturalną $N \geq \frac{4}{\epsilon}$. Następnie rozważmy otwarte sąsiedztwo $V$ orbity $P_0$ w przestrzeni $(X, \rho)$ takie, że $\rho(\textrm{pr}_1(U_i), V) \geq 0$ oraz $d_1(g_x, \textrm{Id}) < \frac{\epsilon}{4}$ dla każdego $x \in V$ (pamiętamy, że $g_x = \textrm{Id}$ dla $x \in P_0$).

$P_0$ jest zbiorem niezmienniczym ze względu na f, na mocy lematu \ref{w_otoczeniu_niezmienniczego_jest_otwarty_zawarty_w_iteracjach} istnieje niepusty zbiór otwarty $W \subseteq V$ taki, że $W \cup f(W) \cup \ldots \cup f^N(W) \subseteq V$. Na mocy lematu \ref{lemat_3_glownego_artykulu_istnieje_orbita_okresowa_krojaca_sie_z_rodzina_otwartych} istnieje punkt okresowy $x_0$ odwzorowania $f$ taki, że $x_0 \in \textrm{pr}_1(U_i)$ oraz orbita $x_0$ kroi się niepusto ze zbiorem $W$. Niech $r > 0$ będzie pierwszą dodatnią liczbą całkowitą dla której $f^r(x_0) \in W$. Wtedy $f^r(x_0), f^{r+1}(x_0), \ldots, f^{r+N-1}(x_0) \in V$. Niech $s \geq 0$ będzie pierwszą nieujemną liczbą całkowitą dla której $f^{r+N+s}(x_0) = x_0$, tj. $r+N+s$ jest okresem punktu $x_0$. Weźmy $y_0$ takie, że $(x_0, y_0) \in U_i$. Ponieważ wszystkie odwzorowania włóknowe ($g_x$) odwzorowania $F$ są ,,na'', to istnieje punkt $y^* \in (0,1)$ taki, że $F^s(f^{r+N}(x_0), y^*) = (x_0, y_0)$ (patrz lemat \ref{wszystkie_wloknowe_na_implikuje_punkt_okresowy}).

Przypadek 1. $z = \textrm{pr}_2(F^r(x_0, y_0))$ jest różne od 0 i 1.
Oznaczmy przez $g$ odwzorowanie z $C(I)$ posiadające następujące trzy własności:
\begin{enumerate}[label=(g\arabic*)]
\item \label{g_jeden} $d_1(g, \mathrm{Id}) < \frac{\epsilon}{4}$,
\item \label{g_dwa} $g$ jest odwzorowaniem niemalejącym, pozostawiającym końce przedziału $I$ niezmienione,
\item \label{g_trzy} $g^N(z) = y^*$.
\end{enumerate}

Następnie, rozważmy odwzorowanie $h \in C(I)$ posiadające trzy następujące własności:
\begin{enumerate}[label=(h\arabic*)]
\item \label{h_jeden} $d_1(h, g_{x_0}) < \frac{\epsilon}{4}$,
\item \label{h_dwa} $h(y_0) = g_{x_0}(y_0)$,
\item \label{h_trzy} $h$ jest stałe na zwartym odcinku $[a,b] \subseteq \textrm{pr}_2(U_i)$ zawierającym punkt $y_0$ w swoim wnętrzu.
\end{enumerate}

Weźmy teraz odwzorowanie $G = (f, \widetilde{g}_x) \in \mathcal{F}$ takie, że $d_2(G, F) < \frac{\epsilon}{2}$ oraz

\begin{equation}
    \widetilde{g}_x =
    \begin{cases}
        h & \text{if $x=x_0$,}\\
        g & \text{if $x \in \{f^k(x_0) : r \leq k \leq r+N-1\}$,}\\
        g_x & \text{if $x \in \{f^k(x_0) : 1 \leq k \leq r-1$ or $r+N \leq k \leq r+N+s-1\}$.}
    \end{cases}
\end{equation}

Odwzorowanie takie istnieje na mocy lematu \ref{lemat_4_glownego_artykulu}, ponieważ
$$d_1(\widetilde{g}_{x_0}, g_{x_0}) = d_1(h, g_{x_0}) < \frac{\epsilon}{4}$$
oraz dla $x \in \{f^r(x_0), f^{r+1}(x_0), \ldots, f^{r+N-1}(x_0)\}$,
$$d_1(\widetilde{g}_x, g_x) = d_1(g, g_x) \leq d_1(g, \mathrm{Id}) + d_1(\mathrm{Id}, g_x) < \frac{\epsilon}{4} + \frac{\epsilon}{4} = \frac{\epsilon}{2}$$.

Punkt $(x_0, y_0) \in U_i$ jest punktem okresowym odwzorowania $G$, ponieważ

\begin{equation}
\begin{split}
G^{r+N+s}(x_0, y_0) & = G^{r+N+s-1}(G(x_0, y_0)) \\
& = G^{r+N+s-1}(F(x_0, y_0)) = G^{N+s}(F^r(x_0, y_0)) \\
& = G^{N+s}(f^r(x_0), z) = G^s(f^{r+N}(x_0), y^*) \\
& = F^s(f^{r+N}(x_0), y^*) = (x_0, y_0).
\end{split}
\end{equation}
\textcolor{red}{TODO: Szczegółowe uzasadnienie powyższej równości:}

Ponadto, ze względu na \ref{h_trzy}, na mocy \ref{obraz_x0_razy_ab_przez_iterate_G_rowna_sie_x0y0} zachodzi
$$G^{r+N+s}(\{x_0\} \times [a, b]) = \{(x_0, y_0)\}.$$

Z faktu $y_0 \in (a, b)$ oraz \ref{obraz_przez_iterate_kazdego_dostatecznie_bliskiego_G_zawiera_sie_w_x0_razy_ab} wynika, że każde odwzorowanie $\widetilde{G} \in \mathcal{F}$ dostatecznie bliskie $G$ posiada własność
$$\widetilde{G}^{r+N+s}(\{x_0\} \times [a, b]) \subseteq \{x_0\} \times [a, b].$$


Zatem $\widetilde{G}$ posiada punkt okresowy w $\{x_0\} \times [a, b] \subseteq U_i$. Zatem $G \in \mathcal{F}_{SO}^i$, co kończy dowód dla przypadku 1.

Przypadek 2. $\textrm{pr}_2(F^r(x_0, y_0))$ jest równy $0$ lub $1$.

W takim przypadku użyjemy lematu \ref{lemat_4_glownego_artykulu} aby dostać odwzorowanie $H = (f, h_x) \in \mathcal{F}$ takie, że $d_2(H, F) < \frac{\epsilon}{2}$ i dla $x \in \{x_0, f(x_0), \ldots, f^{r-1}(x_0)\}$ odwzorowania włóknowe $h_x$ są ściśle rosnące. (Odwzorowanie takie istnieje na mocy \ref{na_mocy_lematu_4_mozemy_dostac_odwzorowanie_H_z_wloknami_scisle_rosnacymi}) Ponieważ $y_0$ jest różnie od $0$ i $1$ dostajemy, że $\textrm{pr}_2(H^r(x_0, y_0))$ również jest różne od $0$ i $1$ (uzasadnienie w \ref{y0_rozne_od_0_i_1_implikuje_pr2Hrx0y0_rozne_od_0_i_1}). Następnie korzystając z przypadku 1 i z lematu \ref{przypadek1_daje_odwzorowanie_G_bliskie_H} dostajemy odwzorowanie $G \in \mathcal{F}_{SO}^i$, dla którego zachodzi nierówność $d_2(G, H) < \frac{\epsilon}{2}$. Wówczas $d_2(G, F) < \epsilon$. (TODO powolac sie na nierownosc trojkata)


\end{proof}




% END MOJE ====================================================



{\backmatter \chapter{Podsumowanie}}
Podsumowanie w pracach matematycznych nie jest obligatoryjne. Warto jednak na zakończenie krótko napisać, co udało nam się zrobić w pracy, a czasem także o tym, czego nie udało się zrobić.

{\backmatter \chapter{Dodatek}}
Dodatek w pracach matematycznych również nie jest wymagany. Można w nim przedstawić np. jakiś dłuższy dowód, który z pewnych przyczyn pominęliśmy we właściwej części pracy lub (np. w przypadku prac statystycznych) umieścić dane, które analizowaliśmy.

%%%%%%%%%%%%%%%%%%%%%%%%%%%%%%%%%%%%%%%%%%%%%%%%%%%%%%%%%
% BIBLIOGRAFIA
% W tworzeniu bibliografii najlepiej korzystać z BibTex'a, 
% który jest częścią systemu Tex. W naszym przypadku funkcję 
% przechowalni literatury, do której się odwołujemy, pełni 
% plik bibliografia.bib. Nie musimy ręcznie dodawać nowych 
% pozycji do bibliografii. Możemy wejść np. na stronę 
% https://mathscinet.ams.org/mathscinet/index.html, 
% znaleźć odpowiednią pozycję, wybrać ją, a następnie zmienić 
% 'Select alternative format' na BibTeX, skopiować uzyskany 
% tekst, wkleić do pliku bibliografia.bib i skompilować. 
% Gotowe informacje do pliku bibliografia.bib można znaleźć 
% także na https://arxiv.org - gdy znajdziemy interesującą nas 
% pracę, szukamy 'References & Citations' i klikamy 'NASA ADS', 
% a potem 'Bibtex entry for this abstract' 
% i postępujemy tak jak wcześniej.
%%%%%%%%%%%%%%%%%%%%%%%%%%%%%%%%%%%%%%%%%%%%%%%%%%%%%%%%%
\newpage
% w nawiasie klamrowym wpisujemy nazwę pliku z bibliografią w formacie .bib
\bibliografia{bibliografia.bib} 
\end{document}